% !TEX root = ../main.tex
% !TEX program = xelatex
\section{模型评价与推广}

\subsection{模型优点与创新}

\begin{enumerate}
    \item \textbf{四问共用一套可以核对的电量流骨架。}
    本文始终用十分钟电量连接负载、光伏、外网和电池，并统一处理90\%充电效率、90\%放电效率、储电上下限及功率限制。各问只在信息条件和费用规则发生变化时增加相应模块，避免同一物理过程被重复解释成不同含义。问题一的确定性调度由线性规划完成[1]，后续各问沿用同一基础关系，便于逐段复算。

    \item \textbf{把预测和高价缺口风险分开处理。}
    负载与光伏分别预测，使星期用电规律和近期光伏变化都能被解释[5]；误差准备则专门处理五倍紧急电价造成的损失不对称。
    问题三进一步用当天已经发生的负载修正后续判断，修正幅度由历史上偏差的延续程度确定。
    分位检验表明，准备不足和准备过多都会增加总费用。因此，模型同时考虑普通订电与缺口补购，不以单一拟合指标或最低紧急电量代替整体经济性。

    \item \textbf{把日内信息的价值拆成两个可说明的部分。}
    在相同模型下，问题三分别比较零点固定计划、日内反馈与改单、加入后发官方预报三种信息条件。
    日内反馈与改单节省50.63万元，新预报额外节省6.95万元。
    这一比较说明两类信息各有什么作用，也直接回答了题目关于日内预报是否值得引入的问题。

    \item \textbf{根据决策方式选择价格模型。}
    问题四没有机械地给两条分支使用同一种算法。零点分支采用提升树表达非线性关系[3]；日内分支采用岭回归稳定处理高度相关的历史价格特征[2]。后者以极小费用差换取更少紧急购电和更清楚的系数含义，体现了经济性、风险和可解释性的综合选择。

    \item \textbf{结论能够追溯到实际运行过程。}
    本文使用SciPy和scikit-learn完成数值计算与候选预测[6,7]，线性规划由HiGHS求解[4]。
    各时段订单与实际供需可以逐项核对，费用能够由普通订单、调整和紧急购电重新计算。
    指定日期分析还区分了库存不足与放电功率不足两种缺口，解释电池仍有余电时为何也可能需要补购，使费用结论与实际运行相互印证。
\end{enumerate}


\subsection{模型局限与改进方向}

首先，历史窗口、误差分位、日期衰减和价格模型参数主要在2025年数据内比较。现有结果能够说明本题数据范围内的取舍，但不能证明同一组参数在其他年份或其他地区仍然最优。若获得跨年数据，可按年份滚动回测，并在不同季节分别校准参数，以检验真正的跨年泛化能力。

其次，本文先预测负载、光伏和价格，再把预测结果送入优化。这样的分层结构容易解释，也便于定位误差来源，但没有同时描述三类不确定量可能共同出现的极端情景。后续可构造少量有现实含义的联合情景，比较连续阴天、高负载和高电价叠加时的计划，进一步控制尾部风险。

再次，题目没有提供电池老化费用、外送电价、网络损耗和可中断负载信息，因此本文未把这些因素写入目标函数。若用于真实微网，可增加循环寿命成本、需求响应和售电约束；末端库存的参考价值也应根据下一日市场规则重新估计，而不能直接视为可出售资产。


\subsection{模型推广性评估}

模型的核心输入只有分时负载、本地可再生能源、设备参数、交易价格和允许的决策时点。替换这些数据并重新校准预测与准备参数后，可用于社区、校园、工业园区及偏远地区的微网运行，也可扩展到多台储能、风光互补、充电站和可调负载。滚动求解方式允许系统随着新数据修正未来计划，适合部署在需要持续更新的能源管理平台中。

从社会应用看，该方法优先消纳本地光伏，并通过储能和合理购电减少高峰期对外网的额外压力，有助于提高可再生能源利用率、降低能源浪费，服务碳达峰和碳中和目标[8,9]。推广时还应把安全供电和居民基本用电放在成本优化之前，使节费收益建立在可靠、公平的能源服务之上。这样的应用方向既符合绿色发展和资源节约要求，也体现以人民为中心、共建共享的发展取向。
